% Citation review: technical citations and bibliography updated; see SOURCE_REVIEW_NOTES.txt for scope and unresolved wording.
\documentclass[11pt,a4paper]{article}
\usepackage[T1]{fontenc}
\usepackage{lmodern,microtype}
\usepackage[margin=20mm,headheight=14pt]{geometry}
\usepackage{amsmath,amssymb,graphicx,booktabs,array,caption,xcolor,fancyhdr}
\usepackage[colorlinks=true,linkcolor=black,citecolor=black,urlcolor=blue]{hyperref}
\usepackage{xurl}
\graphicspath{{outputs/}{./}}
\setlength{\parindent}{0pt}
\setlength{\parskip}{5pt}
\captionsetup{font=small,labelfont=bf,hypcap=false}
\pagestyle{fancy}
\fancyhf{}
\fancyhead[L]{\small MATH 300 \(\cdot\) Video compression with the SVD}
\fancyhead[R]{\small Review draft}
\fancyfoot[C]{\thepage}
\newcommand{\norm}[1]{\left\lVert #1\right\rVert}
\newcommand{\code}[1]{\texttt{#1}}
\begin{document}

\begingroup\centering
{\Large\bfseries Compressing a Push-Up Video with the Incremental SVD}\\[5pt]
Hakim Nizam \(\cdot\) MATH 300, Fall 2026 \(\cdot\) Revised November 20, 2026
\par\endgroup

\begingroup\small
\textbf{Draft status.} This is a final review draft based on measured
results. The methods draw on class materials and cited references.
Contributions and software assistance are summarized on
page~\pageref{sec:class-notes}.
\endgroup

\section*{1. Question and data}

Consecutive frames share structure that a smaller set of basis images can
represent. This project tests whether block incremental SVD matches batch
SVD when all components are kept, and how limiting the retained rank affects
accuracy. It follows Approach A in the brief \cite{brief} and the
column-update method of K\"uhl et al.\ \cite{kuhl}.

I provided the personal push-up recording on November~13, 2026
\cite{video}. The camera is approximately fixed while the body moves.
The portrait source is \(2160\times3840\) pixels. ImageIO's FFmpeg plugin
decoded 184 frames, although the metadata listed 194 \cite{imageio};
the analysis uses the decoded frames.

Keeping every other decoded frame, starting at index zero, gives \(T=92\).
Pillow converts each to grayscale and resizes it to 54 pixels wide by
96 high, preserving the aspect ratio \cite{pillow}. Flattening each image
in the same order \cite{numpy,sep18} gives one column of \(N=5184\) pixel values:
\[
Y=[\,y_1\ y_2\ \cdots\ y_{92}\,]\in\mathbb{R}^{5184\times92}.
\]
Column \(j\) is one frame; row \(i\) tracks one prepared pixel across time.
The pixels are unsigned 8-bit values, while factorizations use floating-point
arithmetic. The mean image is not subtracted.

\section*{2. Batch SVD as a reference}

Batch SVD processes all 92 frame columns
together. Its reduced factorization is
\[
Y=U\Sigma W,\qquad W=V^{\mathsf T}.
\]
NumPy returns \(U\), a vector \code{s}, and \(W\) (called \code{vt} in the
code), with shapes \(5184\times92\), \(92\), and \(92\times92\)
\cite{numpy,nov9}. Each column \(u_i\) of \(U\) is a shared pixel pattern;
reshaping it to \(96\times54\) gives a \emph{basis image}, computed from the
video rather than selected as an original frame. The columns of \(U\) and
rows of \(W\) are orthonormal: unit length and mutually perpendicular within
each set.

\textbf{Patterns and weights.}
\(\Sigma=\operatorname{diag}(s_1,\ldots,s_{92})\) holds nonnegative singular
values, largest first. They measure the overall strength of the corresponding
patterns. For frame \(j\), the actual weight on basis image \(i\) is
\(s_iW_{ij}\). Thus a frame is rebuilt by multiplying each retained basis
image by its own weight and adding the pixel values \cite{numpy,damle}.

\textbf{What rank changes.} A rank-\(q\) approximation keeps the first \(q\)
components and reconstructs all 92 frames:
\begin{equation}
Y_q=\underbrace{U_q}_{5184\times q}
    \underbrace{(\Sigma_qW_q)}_{q\times92}
    \ \in\mathbb{R}^{5184\times92}.
\label{eq:rankq}
\end{equation}
Here \(\Sigma_q\) is \(q\times q\), and \(W_q\) is \(q\times92\).
At rank 5, every frame uses the same five basis images with its own five
weights. The output still has 5184 pixel rows and 92 frame columns.

This is the best approximation of rank at most \(q\) in the Frobenius norm,
which is the square root of the sum of squared matrix entries. Its relative
error, over every pixel of every prepared frame, is \cite{nov13,damle}
\begin{equation}
e_q=\frac{\norm{Y-Y_q}_F}{\norm{Y}_F}
    =\sqrt{\frac{\sum_{i>q}s_i^2}{\sum_i s_i^2}}.
\label{eq:tail}
\end{equation}
Small discarded singular values therefore mean small relative error when
their squared sum is small compared with the total. This formula independently
checks the measured batch error.

\newpage
\section*{3. Incremental update and numerical stability}

\textbf{Frames per update and patterns retained are different choices.}
Incremental SVD updates existing factors as new frame columns arrive.
The block size \(b=10\) means up to ten new frames per update; the final
block contains two. The retained-rank limit \(q\) controls how many basis
images are kept. A rank-20 run therefore still processes ten new frames at
a time. The function starts with an SVD of the first block, then repeats
the following update \cite{kuhl}.

\textbf{1. Find what the existing basis already explains.}
Suppose \(t\) frames are represented by \(\widetilde Y=U\Sigma W\), with
\(r\) currently stored basis images. A new block \(B\) has shape
\(5184\times b\). Compute
\[
M=U^{\mathsf T}B,\qquad P=B-UM,\qquad P=Q_P R_P.
\]
The \(r\times b\) matrix \(M\) gives the new frames' weights in the existing
basis. Thus \(UM\) is the explained part, and the residual \(P\) is what
remains. Reduced QR writes that residual using perpendicular directions
\(Q_P\) and coefficients \(R_P\) \cite{nov6,overleaf}.

\textbf{2. Assemble a smaller update problem.}
Combine the existing basis and residual directions in \(C\), then form
\begin{equation}
C=[\,U\ Q_P\,],\qquad
K=\begin{bmatrix}\Sigma&M\\0&R_P\end{bmatrix},\qquad
D=\begin{bmatrix}W&0\\0&I_b\end{bmatrix}.
\label{eq:core}
\end{equation}
Here \(K\) combines the old singular values with the new-frame coefficients.
For these tall frame blocks, residual QR returns \(b\) columns, so \(K\) is
\((r+b)\times(r+b)\). These factors satisfy
\begin{equation}
[\,\widetilde Y\ B\,]=CKD.
\label{eq:identity}
\end{equation}
The matrix \(D\) carries forward the old temporal coefficients and inserts
the new frames. This identity represents the old approximation plus the new
data in smaller coordinates \cite{kuhl,brand}.

\textbf{3. Restore a perpendicular basis and compute the small SVD.}
Similar frames can make \(P\) small or rank deficient. Some QR directions
may then overlap the old basis. A second QR makes the expanded basis
orthonormal, also addressing roundoff drift \cite{zhang}:
\[
C=Q_C R_C,\qquad
H=R_CK=\widehat U\,\widehat\Sigma\,\widehat W.
\]
\(R_C\) transfers the change of coordinates into \(H\), preserving the
represented pixels: \(CK=Q_CH\). Omitting \(R_C\) would change the matrix.
NumPy computes both QR factorizations and the SVD of \(H\) \cite{numpy}.

\textbf{4. Update the basis images and all frame coefficients.}
The new factors are
\begin{align}
U_{\mathrm{new}}&=Q_C\widehat U,&
\Sigma_{\mathrm{new}}&=\widehat\Sigma,&
W_{\mathrm{new}}&=\widehat W D.
\label{eq:update}
\end{align}
The right factor keeps old and new frame columns together:
\[
W_{\mathrm{new}}
=\bigl[\,\widehat W_{:,1:r}W\quad
        \widehat W_{:,r+1:r+b}\,\bigr].
\]
The split uses the old value of \(r\). Now \(U_{\rm new}\) contains updated
basis images, and \(\Sigma_{\rm new}W_{\rm new}\) contains weights for all
\(t+b\) frames. In \([\,\widetilde Y\ B\,]=Q_CHD\), \(Q_C\) has orthonormal
columns and \(D\) has orthonormal rows. Hence the small SVD preserves the
update's nonzero singular values. Without any truncation, this agrees with
batch SVD up to rounding.

\textbf{5. Optionally truncate, then repeat.}
The function \code{incremental\_svd} in \path{practice.py} keeps at most \(q\)
components after initialization and every update, shortening all three
factors together. A rank-20 run starts with only ten available modes. Once
\(r=20\), a ten-frame update uses a \(30\times30\) core; the last two-frame
update uses a \(22\times22\) core. Truncation is optimal for the retained
approximation plus the new block, but earlier discarded information cannot
be recovered. Therefore it need not give the globally best rank-\(q\)
approximation of the original video.

\newpage
\section*{4. Validation and the full-video result}

\textbf{Known answer.}
The small test matrix is
\[
Z=\begin{bmatrix}
3&0&4&0&0\\
0&4&0&3&0\\
0&0&0&0&2\\
0&0&0&0&0
\end{bmatrix},
\qquad
ZZ^{\mathsf T}=\operatorname{diag}(25,25,4,0).
\]
Its singular values are therefore \(5,5,2,0\). With blocks of two columns,
the second block lies in the existing span and the last block has one
column. The untruncated implementation returns the expected spectrum and
reconstructs \(Z\). Additional checks with block sizes 1, 2, 3 and 5 also
pass. Singular vectors are not compared entry by entry because signs and
bases for repeated singular values are not unique.

\textbf{Full video.}
The batch and untruncated incremental computations both use the same
\(5184\times92\) matrix. The following dimensionless diagnostics are from a
verified run. Last digits may vary with floating-point execution \cite{numpy,sep25}.

\begin{center}
\renewcommand{\arraystretch}{1.4}
\begin{tabular}{@{}lr@{}}
\toprule
Check & Measured value\\
\midrule
\(\norm{Y-U_{\rm inc}\Sigma_{\rm inc}W_{\rm inc}}_F/\norm{Y}_F\)
& \(8.051\times10^{-15}\)\\
\(\norm{s_{\rm inc}-s_{\rm batch}}_2/\norm{s_{\rm batch}}_2\)
& \(8.844\times10^{-16}\)\\
\(\norm{U_{\rm inc}^{\mathsf T}U_{\rm inc}-I}_F\)
& \(1.559\times10^{-14}\)\\
\bottomrule
\end{tabular}
\end{center}

These checks test the reconstructed pixels, singular values, and final
basis orthogonality separately. All are small enough to support agreement
within numerical rounding for this experiment. The orthogonality value is
a final check, rather than a recorded history after every block.

\begin{figure}[h]
\centering
\begin{minipage}{0.49\textwidth}
\includegraphics[width=\linewidth]{report_spectrum.png}
\end{minipage}\hfill
\begin{minipage}{0.49\textwidth}
\includegraphics[width=\linewidth]{report_error.png}
\end{minipage}
\caption{Left: the 92 singular values of the prepared video, with a logarithmic
vertical axis. The batch and untruncated incremental curves overlap.
Right: whole-video relative Frobenius errors for independently computed
rank-1, rank-5 and rank-20 approximations. Lines connect the three tested ranks.
Plots were generated from this experiment using Matplotlib \cite{matplotlib}.}
\label{fig:curves}
\end{figure}

The leading singular values fall quickly, followed by a longer tail
(Figure~\ref{fig:curves}, left). The first mode accounts for about
\(92.60\%\) of \(\norm{Y}_F^2\). Since the frames are uncentered, this includes
the common brightness and background as well as the moving subject.
It is not a measure of how accurately the push-up motion alone is recovered.

\newpage
\section*{5. Rank comparisons and visible detail}

The rank comparisons are an extension beyond the report's required
untruncated check and also prepare evidence for the final poster
\cite{brief}. Each rank uses a separate incremental pass. Table~\ref{tab:errors}
reports errors over all 92 prepared frames, calculated before any display
clipping. The batch values also agree with equation~\eqref{eq:tail}.

\begin{table}[h]
\centering
\begin{tabular}{@{}rrrr@{}}
\toprule
Rank \(q\) & Batch error & Incremental error & Extra error (percentage points)\\
\midrule
1  & \(27.2092\%\) & \(27.2636\%\) & \(0.0544\)\\
5  & \(9.6202\%\)  & \(9.6810\%\)  & \(0.0608\)\\
20 & \(4.1284\%\)  & \(4.1651\%\)  & \(0.0367\)\\
\bottomrule
\end{tabular}
\caption{Relative Frobenius errors for the whole prepared video, with block
size 10. The last column is incremental minus batch error, expressed in
percentage points. Exact values are saved in \texttt{rank\_comparison.csv}.}
\label{tab:errors}
\end{table}

\begin{figure}[h]
\centering
\includegraphics[width=0.98\textwidth]{report_frames.png}
\caption{Sampled frame 47 of 92 from the personal recording \cite{video},
shown with its incremental reconstructions. All panels use the
same \(96\times54\) image shape and grayscale range 0--255. The error values
in Table~\ref{tab:errors} concern the entire video, rather than this frame alone.
The complete batch/incremental image comparison is supplied separately in
\texttt{frame\_reconstructions.png}. This comparison follows the rank-reconstruction idea in the November~13
class material \cite{nov13}.}
\label{fig:frames}
\end{figure}

For batch rank 1, equation~\eqref{eq:tail} gives
\(e_1=\sqrt{1-0.925966}\approx0.2721\).
This connects the captured squared intensity to the measured relative error.
At rank 1, the fixed room is recognizable, but the arm and body positions are
smeared together. A large captured-energy percentage therefore does not
guarantee a clear moving subject. Rank 5 removes much of the visible ghosting,
and rank 20 brings back more of the edge detail in the selected frame
(Figure~\ref{fig:frames}).

Incremental error is slightly higher at every tested rank. Repeated
truncation provides a reason for this gap: a direction that seems weak in an
early block can be useful later. The small gap in this particular clip
supports the usefulness of the update, but does not establish the same
performance for every scene or frame order. The batch approximation remains
the reference optimum at each rank.

The results concern one short fixed-camera clip, one block size, and three
truncation ranks. Grayscale conversion, resizing and temporal sampling
already remove information before SVD. The reported errors measure the
additional approximation of the prepared matrix, and weight stationary
pixels along with motion.

\newpage
\section*{6. Storage, limitations and conclusion}

With singular values stored as a vector, rank-\(q\) factors require
\(q(N+T+1)=5277q\) numbers \cite{nov13}. Table~\ref{tab:storage} uses float64
for the factors. The dense float64 matrix occupies \(3{,}815{,}424\) bytes,
while the prepared uint8 stack occupies \(476{,}928\) bytes.
\begin{center}
\begin{tabular}{@{}rrr@{}}
\toprule
\(q\) & Factor bytes & Dense float64 bytes / factor bytes\\
\midrule
1 & \(42{,}216\) & \(90.38\)\\
5 & \(211{,}080\) & \(18.08\)\\
20 & \(844{,}320\) & \(4.52\)\\
\bottomrule
\end{tabular}
\captionof{table}{Factor-storage estimates with a consistent float64
comparison. At rank 20, the factors exceed the size of the uint8 stack.}
\label{tab:storage}
\end{center}

The script preloads \(Y\), retains a frame list, and makes full reconstructions
for validation. These storage estimates are therefore not measured peak-RAM
savings. A streaming implementation could retain a spatial basis and an
incoming block while updating temporal weights, but storing all temporal
weights still grows with the number of frames \cite{kuhl}. This experiment
does not compare encoded MOV sizes or establish a speed advantage.

Within the scope of the prepared video, the untruncated update
reproduces batch SVD to rounding, and rank 20 gives about \(4.17\%\)
incremental relative error.

\begingroup\small
\phantomsection\label{sec:class-notes}
\textbf{Contributions and AI use.}
I am the sole student participant. I provided the recording, entered and ran
early examples, reported outputs, and reviewed the work. I used GitHub Copilot
in Visual Studio Code to help complete the \(U\)-variable definitions and
matrix setup. The cited class materials, research papers, and software
documentation support the methods and their explanation.
\endgroup

\begingroup\small
\setlength{\parskip}{0pt}
\begin{thebibliography}{16}
\setlength{\itemsep}{0pt}
\bibitem{brief}
S.~Sayyad. \emph{Compressing a Video with the (Incremental) SVD}.
MATH 300 project brief, Fall 2026.
\href{https://github.com/say-yas/wsu_math300/blob/2242bf2f92da2d12b0ac8d933342b008428daaa1/project_proposals/incremental_svd_project.tex}{Official project source}.

\bibitem{video}
H.~Nizam. Personal push-up video, November 13, 2026.
\path{20261113_222244000_iOS.MOV}, provided through OneDrive.

\bibitem{kuhl}
N.~K\"uhl, H.~Fischer, M.~Hinze and T.~Rung (2023).
\emph{An Incremental Singular Value Decomposition Approach for Large-Scale
Spatially Parallel \& Distributed but Temporally Serial Data---Applied to
Technical Flows}. \href{https://arxiv.org/abs/2302.09149v1}{arXiv:2302.09149v1},
Section 2.2, equations 16--19.

\bibitem{brand}
M.~Brand (2006). Fast Low-Rank Modifications of the Thin Singular Value
Decomposition. \emph{Linear Algebra and its Applications} 415, 20--30.
\href{https://www.merl.com/publications/TR2006-059}{MERL TR2006-059}.

\bibitem{numpy}
NumPy. Version 2.4.6:
\href{https://numpy.org/doc/2.4/reference/generated/numpy.linalg.svd.html}{SVD},
\href{https://numpy.org/doc/2.4/reference/generated/numpy.linalg.qr.html}{QR},
\href{https://numpy.org/doc/2.4/reference/generated/numpy.ravel.html}{ravel},
\href{https://numpy.org/doc/2.4/reference/generated/numpy.column_stack.html}{stacking},
\href{https://numpy.org/doc/2.4/user/basics.broadcasting.html}{broadcasting},
\href{https://numpy.org/doc/2.4/reference/generated/numpy.linalg.norm.html}{norms},
\href{https://numpy.org/doc/2.4/reference/generated/numpy.allclose.html}{allclose}.

\bibitem{nov6}
S.~Sayyad. MATH 300, November 6, 2026.
\href{https://github.com/say-yas/wsu_math300/blob/ef5f38f850dabb9a1f19068884ba4b0fa7bb2a5b/Nov6_decompositions/decompositions.ipynb}{Matrix decompositions, Part B (QR)}.

\bibitem{nov9}
S.~Sayyad. MATH 300, November 9, 2026.
\href{https://github.com/say-yas/wsu_math300/blob/435aab704ac2701d771769e232c75295ec97bba9/Nov9_SVD_similarity_transformations/svd-eigenspaces.ipynb}{SVD and eigenspaces, Part A}.

\bibitem{nov13}
S.~Sayyad. MATH 300, November 13, 2026.
\href{https://github.com/say-yas/wsu_math300/blob/2242bf2f92da2d12b0ac8d933342b008428daaa1/Nov13_Imagecompression/svd_image_compression.ipynb}{Image compression with the SVD}.

\bibitem{imageio}
ImageIO. \href{https://imageio.readthedocs.io/en/stable/_autosummary/imageio.v3.imiter.html}{\code{imageio.v3.imiter}}
and \href{https://imageio.readthedocs.io/en/stable/_autosummary/imageio.plugins.ffmpeg.html}{FFmpeg plugin}
documentation. Experiment version 2.37.3.

\bibitem{pillow}
Pillow. \href{https://pillow.readthedocs.io/en/stable/reference/Image.html}{\code{Image.convert} and \code{Image.resize}}
documentation. Experiment version 12.3.0.

\bibitem{matplotlib}
Matplotlib. \href{https://matplotlib.org/stable/api/_as_gen/matplotlib.pyplot.semilogy.html}{\code{pyplot.semilogy}}
and \href{https://matplotlib.org/stable/api/_as_gen/matplotlib.axes.Axes.imshow.html}{\code{Axes.imshow}}
documentation. Experiment version 3.11.0.

\bibitem{overleaf}
S.~Sayyad. \href{https://www.overleaf.com/project/6a8b1c59a232e8dbbb5b0923}{WSU\_MATH300\_Shared},
Activities XXX, XXXI and XXXII. Checked November 19, 2026.

\bibitem{sep18}
S.~Sayyad. MATH 300, September 18, 2026.
\href{https://github.com/say-yas/wsu_math300/blob/8acbae688a5e5ef588dca37c5ca6f17b00600ac4/Sep18/class_intro_numpy_2.ipynb}{NumPy arrays, shapes, reshaping and matrix products}.

\bibitem{sep25}
S.~Sayyad. MATH 300, September 25, 2026.
\href{https://github.com/say-yas/wsu_math300/blob/2e41141c9a3646473619b922458d42a1e5333f66/Sep25/complexity_fp_convergence_practice.ipynb}{Floating-point arithmetic, tolerances and relative error}.

\bibitem{damle}
A.~Damle. \href{https://www.cs.cornell.edu/courses/cs3220/2020fa/SVD.pdf}{CS 3220: The Singular Value Decomposition}.
Cornell University, Fall 2020 course (notes undated), pp.~2--3, Theorem~1.

\bibitem{zhang}
Y.~Zhang (2022). \emph{An answer to an open question in the incremental SVD}.
\href{https://arxiv.org/abs/2204.05398v2}{arXiv:2204.05398v2}, Section~2, pp.~3--6.
\end{thebibliography}
\endgroup
\end{document}

